\subsubsection{Ausblick: Komplexe Zahlen}
\label{subsubsec:fundamentals_complex_numbers}

Um dennoch die Wurzel aus $-1$ ziehen zu können, führt man die imaginäre
Einheit $i$ mit der Eigenschaft
\[
  i^2=-1
\]
ein. Da $i\notin\R$ gilt, ist $i$ keine reelle Zahl. Auf diese Weise erweitert
man die reellen Zahlen $\R$ zu den sogenannten \emph{komplexen Zahlen} $\C$.

\begin{center}
  \begingroup
    \def\NumberSetStage{5}
    \CachedTikzFigure{content/resources/chapters/fundamentals/tikz/number_set_inclusions.tex}
  \endgroup
\end{center}

Eine ausführliche Einführung in die komplexen Zahlen $\C$ folgt in
Kapitel~\ref{chap:complex_numbers}.

